% ------------------------------------------------------------------------
% file `UAA4-exercise-6-solution-body.tex'
%   in folder `xsim/'
%
%     solution of type `exercise' with id `6'
%
% generated by the `solution' environment of the
%   `xsim' package v0.21 (2022/02/12)
% from source `UAA4' on 2022/05/10 on line 77
% ------------------------------------------------------------------------
    \begin{enumerate}
        \item $CB \sslash ED \Rightarrow
                  \begin{aligned}[t]
                      \dfrac{x}{22} & = \dfrac{7}{6}          \\
                      x             & = 22 \cdot \dfrac{7}{6} \\
                      x             & = \dfrac{77}{3}
                  \end{aligned}$
        \item $CD \sslash ED \Rightarrow
                  \begin{aligned}[t]
                      \dfrac{x}{4} & = \dfrac{11}{6}        \\
                      x            & = 4 \cdot\dfrac{11}{6} \\
                      x            & = \dfrac{22}{3}
                  \end{aligned}$
        \item $\dfrac{4,2}{2} = \dfrac{6,3}{3} = \dfrac{8,4}{4} = 2,1$\\
              Les deux triangles ont leur côtés homologues proportionnels (CCC)\\
              $\Rightarrow |\widehat{D}| = \widehat{A}| = 104\degres$

        \item
              \begin{tblr}[t]{hlines, vlines}
                  $ACE$           & $BDE$  & Justifications                 \\
                  $|AE|$          & $|EB|$ & $E$ est le milieu de $[AB]$    \\
                  $|\widehat{E}|$ & $|\widehat{E}|$ & {ils sont alterne internes respectivement avec les angles\\du triangle équiltéral et $AB \sslash CD$.}\\
                  $|CE|$          & $|DE|$ & côtés d'un triangle équiltéral \\
              \end{tblr}
              Par le cas d'isométrie $CAC$, les triangles sont isométriques. Donc $x=6$.
    \end{enumerate}
